%% ODER: format ==         = "\mathrel{==}"
%% ODER: format /=         = "\neq "
%
%
\makeatletter
\@ifundefined{lhs2tex.lhs2tex.sty.read}%
  {\@namedef{lhs2tex.lhs2tex.sty.read}{}%
   \newcommand\SkipToFmtEnd{}%
   \newcommand\EndFmtInput{}%
   \long\def\SkipToFmtEnd#1\EndFmtInput{}%
  }\SkipToFmtEnd

\newcommand\ReadOnlyOnce[1]{\@ifundefined{#1}{\@namedef{#1}{}}\SkipToFmtEnd}
\usepackage{amstext}
\usepackage{amssymb}
\usepackage{stmaryrd}
\DeclareFontFamily{OT1}{cmtex}{}
\DeclareFontShape{OT1}{cmtex}{m}{n}
  {<5><6><7><8>cmtex8
   <9>cmtex9
   <10><10.95><12><14.4><17.28><20.74><24.88>cmtex10}{}
\DeclareFontShape{OT1}{cmtex}{m}{it}
  {<-> ssub * cmtt/m/it}{}
\newcommand{\texfamily}{\fontfamily{cmtex}\selectfont}
\DeclareFontShape{OT1}{cmtt}{bx}{n}
  {<5><6><7><8>cmtt8
   <9>cmbtt9
   <10><10.95><12><14.4><17.28><20.74><24.88>cmbtt10}{}
\DeclareFontShape{OT1}{cmtex}{bx}{n}
  {<-> ssub * cmtt/bx/n}{}
\newcommand{\tex}[1]{\text{\texfamily#1}}	% NEU

\newcommand{\Sp}{\hskip.33334em\relax}


\newcommand{\Conid}[1]{\mathit{#1}}
\newcommand{\Varid}[1]{\mathit{#1}}
\newcommand{\anonymous}{\kern0.06em \vbox{\hrule\@width.5em}}
\newcommand{\plus}{\mathbin{+\!\!\!+}}
\newcommand{\bind}{\mathbin{>\!\!\!>\mkern-6.7mu=}}
\newcommand{\rbind}{\mathbin{=\mkern-6.7mu<\!\!\!<}}% suggested by Neil Mitchell
\newcommand{\sequ}{\mathbin{>\!\!\!>}}
\renewcommand{\leq}{\leqslant}
\renewcommand{\geq}{\geqslant}
\usepackage{polytable}

%mathindent has to be defined
\@ifundefined{mathindent}%
  {\newdimen\mathindent\mathindent\leftmargini}%
  {}%

\def\resethooks{%
  \global\let\SaveRestoreHook\empty
  \global\let\ColumnHook\empty}
\newcommand*{\savecolumns}[1][default]%
  {\g@addto@macro\SaveRestoreHook{\savecolumns[#1]}}
\newcommand*{\restorecolumns}[1][default]%
  {\g@addto@macro\SaveRestoreHook{\restorecolumns[#1]}}
\newcommand*{\aligncolumn}[2]%
  {\g@addto@macro\ColumnHook{\column{#1}{#2}}}

\resethooks

\newcommand{\onelinecommentchars}{\quad-{}- }
\newcommand{\commentbeginchars}{\enskip\{-}
\newcommand{\commentendchars}{-\}\enskip}

\newcommand{\visiblecomments}{%
  \let\onelinecomment=\onelinecommentchars
  \let\commentbegin=\commentbeginchars
  \let\commentend=\commentendchars}

\newcommand{\invisiblecomments}{%
  \let\onelinecomment=\empty
  \let\commentbegin=\empty
  \let\commentend=\empty}

\visiblecomments

\newlength{\blanklineskip}
\setlength{\blanklineskip}{0.66084ex}

\newcommand{\hsindent}[1]{\quad}% default is fixed indentation
\let\hspre\empty
\let\hspost\empty
\newcommand{\NB}{\textbf{NB}}
\newcommand{\Todo}[1]{$\langle$\textbf{To do:}~#1$\rangle$}

\EndFmtInput
\makeatother
%
%
%
%
%
%
% This package provides two environments suitable to take the place
% of hscode, called "plainhscode" and "arrayhscode". 
%
% The plain environment surrounds each code block by vertical space,
% and it uses \abovedisplayskip and \belowdisplayskip to get spacing
% similar to formulas. Note that if these dimensions are changed,
% the spacing around displayed math formulas changes as well.
% All code is indented using \leftskip.
%
% Changed 19.08.2004 to reflect changes in colorcode. Should work with
% CodeGroup.sty.
%
\ReadOnlyOnce{polycode.fmt}%
\makeatletter

\newcommand{\hsnewpar}[1]%
  {{\parskip=0pt\parindent=0pt\par\vskip #1\noindent}}

% can be used, for instance, to redefine the code size, by setting the
% command to \small or something alike
\newcommand{\hscodestyle}{}

% The command \sethscode can be used to switch the code formatting
% behaviour by mapping the hscode environment in the subst directive
% to a new LaTeX environment.

\newcommand{\sethscode}[1]%
  {\expandafter\let\expandafter\hscode\csname #1\endcsname
   \expandafter\let\expandafter\endhscode\csname end#1\endcsname}

% "compatibility" mode restores the non-polycode.fmt layout.

\newenvironment{compathscode}%
  {\par\noindent
   \advance\leftskip\mathindent
   \hscodestyle
   \let\\=\@normalcr
   \let\hspre\(\let\hspost\)%
   \pboxed}%
  {\endpboxed\)%
   \par\noindent
   \ignorespacesafterend}

\newcommand{\compaths}{\sethscode{compathscode}}

% "plain" mode is the proposed default.
% It should now work with \centering.
% This required some changes. The old version
% is still available for reference as oldplainhscode.

\newenvironment{plainhscode}%
  {\hsnewpar\abovedisplayskip
   \advance\leftskip\mathindent
   \hscodestyle
   \let\hspre\(\let\hspost\)%
   \pboxed}%
  {\endpboxed%
   \hsnewpar\belowdisplayskip
   \ignorespacesafterend}

\newenvironment{oldplainhscode}%
  {\hsnewpar\abovedisplayskip
   \advance\leftskip\mathindent
   \hscodestyle
   \let\\=\@normalcr
   \(\pboxed}%
  {\endpboxed\)%
   \hsnewpar\belowdisplayskip
   \ignorespacesafterend}

% Here, we make plainhscode the default environment.

\newcommand{\plainhs}{\sethscode{plainhscode}}
\newcommand{\oldplainhs}{\sethscode{oldplainhscode}}
\plainhs

% The arrayhscode is like plain, but makes use of polytable's
% parray environment which disallows page breaks in code blocks.

\newenvironment{arrayhscode}%
  {\hsnewpar\abovedisplayskip
   \advance\leftskip\mathindent
   \hscodestyle
   \let\\=\@normalcr
   \(\parray}%
  {\endparray\)%
   \hsnewpar\belowdisplayskip
   \ignorespacesafterend}

\newcommand{\arrayhs}{\sethscode{arrayhscode}}

% The mathhscode environment also makes use of polytable's parray 
% environment. It is supposed to be used only inside math mode 
% (I used it to typeset the type rules in my thesis).

\newenvironment{mathhscode}%
  {\parray}{\endparray}

\newcommand{\mathhs}{\sethscode{mathhscode}}

% texths is similar to mathhs, but works in text mode.

\newenvironment{texthscode}%
  {\(\parray}{\endparray\)}

\newcommand{\texths}{\sethscode{texthscode}}

% The framed environment places code in a framed box.

\def\codeframewidth{\arrayrulewidth}
\RequirePackage{calc}

\newenvironment{framedhscode}%
  {\parskip=\abovedisplayskip\par\noindent
   \hscodestyle
   \arrayrulewidth=\codeframewidth
   \tabular{@{}|p{\linewidth-2\arraycolsep-2\arrayrulewidth-2pt}|@{}}%
   \hline\framedhslinecorrect\\{-1.5ex}%
   \let\endoflinesave=\\
   \let\\=\@normalcr
   \(\pboxed}%
  {\endpboxed\)%
   \framedhslinecorrect\endoflinesave{.5ex}\hline
   \endtabular
   \parskip=\belowdisplayskip\par\noindent
   \ignorespacesafterend}

\newcommand{\framedhslinecorrect}[2]%
  {#1[#2]}

\newcommand{\framedhs}{\sethscode{framedhscode}}

% The inlinehscode environment is an experimental environment
% that can be used to typeset displayed code inline.

\newenvironment{inlinehscode}%
  {\(\def\column##1##2{}%
   \let\>\undefined\let\<\undefined\let\\\undefined
   \newcommand\>[1][]{}\newcommand\<[1][]{}\newcommand\\[1][]{}%
   \def\fromto##1##2##3{##3}%
   \def\nextline{}}{\) }%

\newcommand{\inlinehs}{\sethscode{inlinehscode}}

% The joincode environment is a separate environment that
% can be used to surround and thereby connect multiple code
% blocks.

\newenvironment{joincode}%
  {\let\orighscode=\hscode
   \let\origendhscode=\endhscode
   \def\endhscode{\def\hscode{\endgroup\def\@currenvir{hscode}\\}\begingroup}
   %\let\SaveRestoreHook=\empty
   %\let\ColumnHook=\empty
   %\let\resethooks=\empty
   \orighscode\def\hscode{\endgroup\def\@currenvir{hscode}}}%
  {\origendhscode
   \global\let\hscode=\orighscode
   \global\let\endhscode=\origendhscode}%

\makeatother
\EndFmtInput
%



%% ----------- Applications --------------
%% %format applyx a b = a "@" b
%% %format applyx a b = a "\," b

%% Function application with no arguments

%% Function application with 1 argument

%% Function application with 1 argument, without parenthesis

%% Function application with many arguments

%% ----------- End of Applications --------------


%% Sequences e1; ...; en; v

%% xs etc stole syntax used in examples.

%% Old version: | for BNF, [] for choice
%%%format choice = "[ \! ]"
%%%format `choice` = "\mathop{[ \! ]}"

%% New version: tall | for BNF, fat | for choice
%%format vbar = "\mathop{\bigm|}"
%%format choice = "\boldsymbol{|}"
%%format `choice` = "\mathop{\boldsymbol{|}}"




%% Arrays and tuples
%% %format array (x) = "\textbf{array}\{" x "\}"
%% %format arrKW = "\textbf{arr}"
%% %format arr (x) = arrKW "\{" x "\}"

%% %format defKW = "\textbf{def}"

% only used in comments, but still prettier this way:
%%%%format map = "\textbf{map}"
% format := = "\mapsto"
%  format effs x = "\!\!<\!\!" x "\!\!>\!\!" dots
































%% ----------- Metatheory --------------
%%%%%%%%%%%%%%%%%%%%%%%%%%%%%%%%%%%%

\newcommand{\rulegroupname}[1]{\multicolumn{4}{l}{\textit{#1}}\\}
\newcommand{\largerhs}[1]{\multicolumn{2}{@{}l}{#1}}

\newcommand{\figureRulesI}{
\begin{figure}
\begin{longtable}{@{\quad} l @{\hspace{-.5em}} r @{\;}>{$}c<{$}@{\;} l l}
% The quad indents the lines
% The \hspace{-..} removes redundant space between rule name column and LHS column
% that arises because the longest rule name and the longer LHS are not on the same line
% The \; reduces the spacing around the arrow
% The >{$} turns on math mode for the arrow column

\rulegroupname{Primitive operations}

\rulename{p-add} & \ensuremath{\textbf{add}\langle\Varid{k}_{\mathrm{1}},\Varid{k}_{\mathrm{2}}\rangle} & \movesto & \ensuremath{\Varid{k}_{\mathrm{1}}\mathbin{+}\Varid{k}_{\mathrm{2}}} \\
\rulename{p-gt} & \ensuremath{\textbf{gt}\langle\Varid{k}_{\mathrm{1}},\Varid{k}_{\mathrm{2}}\rangle} & \movesto &
    $\begin{cases}
      \ensuremath{\Varid{k}_{\mathrm{1}}} & \text{if $k_1 > k_2$} \\
      \ensuremath{\textbf{fail}} & \text{otherwise}
      \end{cases}$ \\

\iffullverse
    \rulename{p-isInt} & \ensuremath{\textbf{isInt}(\mathit{hnf})} & \movesto &
    $\left\{\begin{array}{@{}ll}
      \ensuremath{()} & \text{if \ensuremath{\mathit{hnf}} is an integer} \\
      \ensuremath{\textbf{fail}} & \text{otherwise}
      \end{array}\right.$
\fi

\rulegroupname{Application}

\rulename{app-beta} & \ensuremath{(\lambda\Varid{x}.\,\Varid{e})(\Varid{v})} & \movesto & \ensuremath{\exists\Varid{x}.\,\Varid{x}\mathrel{=}\Varid{v};\,\Varid{e}}
          & if $x \not\in \freevars{\ensuremath{\Varid{v}}}$ \\
\rulename{app-tup} & \ensuremath{\langle{}\Varid{v}_{\mathrm{0}},\cdots\;\Varid{v}_{\Varid{n}}\rangle(\Varid{v})} & \movesto & \largerhs{\ensuremath{(\Varid{v}\mathrel{=}\mathrm{0};\,\Varid{v}_{\mathrm{0}})\;\rule[-0.3mm]{0.5mm}{2.5mm}\;\cdots\;\rule[-0.3mm]{0.5mm}{2.5mm}\;(\Varid{v}\mathrel{=}\Varid{n};\,\Varid{v}_{\Varid{n}})}} \\

\\[-2mm]
\rulegroupname{Unification}

\rulename{ulit} & \ensuremath{\Varid{k}_{\mathrm{1}}=\Varid{k}_{\mathrm{2}}} & \movesto
    & \largerhs{$\begin{cases}
        \ensuremath{\Varid{k}_{\mathrm{1}}} & \text{if $k_1=k_2$} \\
        \ensuremath{\textbf{fail}} & \text{otherwise}
      \end{cases}$} \\
\rulename{utup} & \llap{\ensuremath{\langle{}\Varid{v}_{\mathrm{1}},\cdots,\Varid{v}_{\Varid{m}}\rangle=\langle{}\Varid{v}_{1}',\cdots,\Varid{v}_{\Varid{n}}'\rangle}}
    % the \llap allows the LHS to reach into the first columns; ok due to the short rule name
    & \movesto & \largerhs{$\begin{cases}
        \multicolumn{2}{@{}l}{\ensuremath{\Varid{v}_{\mathrm{1}}=\Varid{v}_{1}';\,\cdots;\,\Varid{v}_{\Varid{m}}\mathrel{=}\Varid{v}_{\Varid{m}}';\,\langle{}\Varid{v}_{\mathrm{1}},\cdots,\Varid{v}_{\Varid{m}}\rangle}} \\
            & \text{if $m=n$} \\
        \ensuremath{\textbf{fail}} & \text{otherwise}
      \end{cases}
    $}\\
\rulename{ux1} & \ensuremath{\Varid{k}=\langle{}\Varid{v}_{\mathrm{1}},\cdots,\Varid{v}_{\Varid{n}}\rangle} & \movesto & \ensuremath{\textbf{fail}} \\
\rulename{ux2} & \ensuremath{\langle{}\Varid{v}_{\mathrm{1}},\cdots,\Varid{v}_{\Varid{n}}\rangle=\Varid{k}} & \movesto & \ensuremath{\textbf{fail}} \\
\rulename{ux3} & \ensuremath{(\lambda\Varid{x}.\,\Varid{e})=\mathit{hnf}} & \movesto & \ensuremath{\textbf{fail}}\\
\rulename{ux4} & \ensuremath{\mathit{hnf}=(\lambda\Varid{x}.\,\Varid{e})} & \movesto & \ensuremath{\textbf{fail}}\\
\rulename{ux5} & \ensuremath{\Varid{op}=\mathit{hnf}} & \movesto & \ensuremath{\textbf{fail}}\\
\rulename{ux6} & \ensuremath{\mathit{hnf}=\Varid{op}} & \movesto & \ensuremath{\textbf{fail}}\\
\rulename{ux-occurs} & \ensuremath{\Varid{x}\mathrel{=}\Conid{V} [\mskip1.5mu \Varid{x}\mskip1.5mu]} & \movesto & \ensuremath{\textbf{fail}} & if $\ensuremath{\Conid{V}} \ne \ensuremath{\hole}$ \\

\\[-2mm]
\rulegroupname{Unification variables}

\rulename{subst} & \ensuremath{X [\mskip1.5mu \Varid{x}=\Varid{v}\mskip1.5mu]} & \movesto & \ensuremath{(\subst{X}{\Varid{x}}{\Varid{v}}) [\mskip1.5mu \Varid{x}=\Varid{v}\mskip1.5mu]}
    & if $\ensuremath{\Varid{x}} \in \freevars{\ensuremath{X}}$, $\ensuremath{\Varid{x}} \notin \freevars{v}$ \\
\rulename{subst-rec} & \ensuremath{\Varid{x}\mathrel{=}\Conid{V} [\mskip1.5mu \lambda\Varid{y}.\,\Varid{e}\mskip1.5mu]} & \movesto & \largerhs{\ensuremath{\Varid{x}\mathrel{=}\Conid{V} [\mskip1.5mu \lambda\Varid{y}.\,\exists\Varid{x}.\,\Varid{x}\mathrel{=}\Conid{V} [\mskip1.5mu \lambda\Varid{y}.\,\Varid{e}\mskip1.5mu];\,\Varid{e}\mskip1.5mu]}} \\
    &&&& if $\ensuremath{\Varid{x}} \in \freevars{\ensuremath{\lambda\Varid{y}.\,\Varid{e}}}$ \\
\rulename{def-eliml} & \ensuremath{\exists\Varid{x}\;\Varid{y}_{\mathrm{1}}\;\cdots\;\Varid{y}_{\Varid{n}}.\,X [\mskip1.5mu \Varid{x}=\Varid{v}\mskip1.5mu]} & \movesto & \ensuremath{\exists\Varid{y}_{\mathrm{1}}\;\cdots\;\Varid{y}_{\Varid{n}}.\,X [\mskip1.5mu \Varid{v}\mskip1.5mu]}
    & if $\ensuremath{\Varid{x}} \notin \freevars{\ensuremath{X}}\cup \freevars{\ensuremath{\Varid{v}}}, \ensuremath{\Varid{x}} \neq \ensuremath{\Varid{y}_{\Varid{i}}}$ \\
\rulename{def-elimr} & \ensuremath{\exists\Varid{x}\;\Varid{y}_{\mathrm{1}}\;\cdots\;\Varid{y}_{\Varid{n}}.\,X [\mskip1.5mu \Varid{v}\mathrel{=}\Varid{x}\mskip1.5mu]} & \movesto & \ensuremath{\exists\Varid{y}_{\mathrm{1}}\;\cdots\;\Varid{y}_{\Varid{n}}.\,X [\mskip1.5mu \Varid{v}\mskip1.5mu]}
    & if $\ensuremath{\Varid{x}} \notin \freevars{\ensuremath{X}}\cup \freevars{\ensuremath{\Varid{v}}}, \ensuremath{\Varid{x}} \neq \ensuremath{\Varid{y}_{\Varid{i}}}$ \\
\rulename{def-elimv} & \ensuremath{\exists\Varid{x}\;\Varid{y}_{\mathrm{1}}\;\cdots\;\Varid{y}_{\Varid{n}}.\,X [\mskip1.5mu \Varid{y}\mathrel{=}\Varid{x}\mskip1.5mu]} & \movesto & \ensuremath{\exists\Varid{y}_{\mathrm{1}}\;\cdots\;\Varid{y}_{\Varid{n}}.\,\subst{(X [\mskip1.5mu \Varid{v}\mskip1.5mu])}{\Varid{x}}{\Varid{y}}}
    & if $\ensuremath{\Varid{x}} \neq \ensuremath{\Varid{y}}, \ensuremath{\Varid{x}} \neq \ensuremath{\Varid{y}_{\Varid{i}}}$ \lennart{NEW} \\
\rulename{swap} & \ensuremath{\mathit{hnf}=\Varid{x}}    & \movesto & \ensuremath{\Varid{x}=\mathit{hnf}} \\
\rulename{def-float} & \ensuremath{X [\mskip1.5mu \exists\Varid{x}.\,\Varid{e}\mskip1.5mu]} & \movesto & \ensuremath{\exists\Varid{x}.\,X [\mskip1.5mu \Varid{e}\mskip1.5mu]}
    & if $\ensuremath{X} \ne \hole,\,  x \notin \freevars{\ensuremath{X}}$\\

\\[-2mm]
\rulegroupname{Sequencing}

\rulename{seq} & \ensuremath{\Varid{v};\,\Varid{e}} & \movesto & \ensuremath{\Varid{e}} \\
\rulename{seq-assoc} & \ensuremath{(\Varid{e}_{\mathrm{1}};\,\Varid{e}_{\mathrm{2}});\,\Varid{e}_{\mathrm{3}}} & \movesto & \ensuremath{\Varid{e}_{\mathrm{1}};\,(\Varid{e}_{\mathrm{2}};\,\Varid{e}_{\mathrm{3}})} \\
\rulename{unify-seql} & \ensuremath{(\Varid{e}_{\mathrm{1}};\,\Varid{e}_{\mathrm{2}})=\Varid{e}_{\mathrm{3}}} & \movesto & \ensuremath{\Varid{e}_{\mathrm{1}};\,(\Varid{e}_{\mathrm{2}}=\Varid{e}_{\mathrm{3}})} \\
\rulename{unify-seqr} & \ensuremath{\Varid{v}=(\Varid{e}_{\mathrm{1}};\,\Varid{e}_{\mathrm{2}})} & \movesto & \ensuremath{\Varid{e}_{\mathrm{1}};\,(\Varid{v}=\Varid{e}_{\mathrm{2}})} \\
\rulename{unify-unifyl} & \ensuremath{(\Varid{e}_{\mathrm{1}}=\Varid{e}_{\mathrm{2}})=\Varid{e}_{\mathrm{3}}} & \movesto & \ensuremath{\exists\Varid{x}.\,\Varid{x}\mathrel{=}\Varid{e}_{\mathrm{1}};\,\Varid{x}\mathrel{=}\Varid{e}_{\mathrm{2}};\,\Varid{x}\mathrel{=}\Varid{e}_{\mathrm{3}}} & \ensuremath{\Varid{x}} fresh\\
\rulename{unify-unifyr} & \ensuremath{\Varid{e}_{\mathrm{1}}=(\Varid{e}_{\mathrm{2}}=\Varid{e}_{\mathrm{3}})} & \movesto & \ensuremath{\exists\Varid{x}.\,\Varid{x}\mathrel{=}\Varid{e}_{\mathrm{1}};\,\Varid{x}\mathrel{=}\Varid{e}_{\mathrm{2}};\,\Varid{x}\mathrel{=}\Varid{e}_{\mathrm{3}}} & \ensuremath{\Varid{x}} fresh\\

\\[-2mm]
\rulegroupname{Fail propagation}

\rulename{fail-def} & \ensuremath{\exists\Varid{x}.\,\textbf{fail}} & \movesto & \ensuremath{\textbf{fail}}  \\
\rulename{fail}     & \ensuremath{X [\mskip1.5mu \textbf{fail}\mskip1.5mu]} & \movesto & \ensuremath{\textbf{fail}} & if $\ensuremath{X} \ne \hole$

\end{longtable}

  \caption{The Verse calculus: rewrite rules (1)}
    \label{fig:general-axioms} \label{fig:rules1}
\end{figure}
}

\newcommand{\figureContexts}{
\begin{figure}
  $$
    \begin{array}{lr@{\;}c@{\;}l}
    \text{Execution contexts} & X & ::= & \hole \vb \ensuremath{X=\Varid{e}} \vb \ensuremath{\Varid{e}=X} \vb \ensuremath{X;\,\Varid{e}} \vb \ensuremath{\Varid{e};\,X} \\
    \text{Value contexts} & V & ::= & \hole \vb \ensuremath{\langle{}\Varid{v}_{\mathrm{1}},\cdots,\Conid{V},\cdots,\Varid{v}_{\Varid{n}}\rangle}\\
    \text{Scope contexts} & SX & ::= & \ensuremath{\hole\;\rule[-0.3mm]{0.5mm}{2.5mm}\;\Varid{e}} \vb \ensuremath{\Varid{e}\;\rule[-0.3mm]{0.5mm}{2.5mm}\;\hole}
    \vb \ensuremath{\textbf{one}\{\hole\}} \vb \ensuremath{\textbf{all}\{\hole\}} \\
\iffullverse
    && \vb & \ensuremath{\textbf{succeeds}\,\{\hole\}} \vb \ensuremath{\textbf{decides}\,\{\hole\}} \\
\fi
    \text{Choice contexts} & CX & ::= & \hole \vb \ensuremath{CX=\Varid{e}} \vb \ensuremath{\Varid{ce}=CX} \vb \ensuremath{CX;\,\Varid{e}} \vb \ensuremath{\Varid{ce};\,CX} \vb \ensuremath{\exists\Varid{x}.\,CX} \\
    \text{Choice-free exprs} & ce  & ::=  & \ensuremath{\Varid{v}} \vb \ensuremath{\Varid{ce}_{\mathrm{1}}=\Varid{ce}_{\mathrm{2}}} \vb \ensuremath{\Varid{ce}_{\mathrm{1}};\,\Varid{ce}_{\mathrm{2}}}
    \vb \ensuremath{\textbf{one}\{\Varid{e}\}} \vb \ensuremath{\textbf{all}\{\Varid{e}\}}
    %%% NOTE: the following is not in the POPL submission,
    %%% but needed to make progress in (1,2)[x:int]
    \vb \ensuremath{\Varid{op}(\Varid{v})}
  \end{array}
   $$
  \caption{The syntax of contexts}
    \label{fig:contexts}
\end{figure}
}

\newcommand{\figureRulesII}{
\begin{figure}
  $$
  \begin{array}{l}
    \textit{One and all} \\
    \begin{array}[t]{lr@{\;}c@{\;}l}
    \rulename{one-fail} & \ensuremath{\textbf{one}\{\textbf{fail}\}} & \movesto & \ensuremath{\textbf{fail}} \\
    \rulename{one-choice} & \ensuremath{\textbf{one}\{\Varid{v}\hspace{1mm}\mathop{\rule[-0.3mm]{0.5mm}{2.5mm}}\hspace{1mm}\Varid{e}\}} & \movesto & \ensuremath{\Varid{v}} \\
    \rulename{one-value} & \ensuremath{\textbf{one}\{\Varid{v}\}} & \movesto & \ensuremath{\Varid{v}} \\
    \end{array} \hspace{5mm}
    \begin{array}[t]{lr@{\;}c@{\;}l}
    \rulename{all-fail} & \ensuremath{\textbf{all}\{\textbf{fail}\}} & \movesto & \ensuremath{\langle{}\rangle} \\
    \rulename{all-choice} &
      \ensuremath{\textbf{all}\{\Varid{v}_{\mathrm{1}}\;\rule[-0.3mm]{0.5mm}{2.5mm}\;\cdots\;\rule[-0.3mm]{0.5mm}{2.5mm}\;\Varid{v}_{\Varid{n}}\}} & \movesto & \ensuremath{\langle{}\Varid{v}_{\mathrm{1}},\cdots,\Varid{v}_{\Varid{n}}\rangle} \\
    \rulename{all-value} & \ensuremath{\textbf{all}\{\Varid{v}\}} & \movesto & \ensuremath{\langle{}\Varid{v}\rangle} \\
    \end{array} \\

    \\[-2mm]

    \textit{Choice} \\
    \begin{array}{lr@{\;}c@{\;}ll}
    \rulename{fail-l} & \ensuremath{\Conid{SX} [\mskip1.5mu \textbf{fail}\;\rule[-0.3mm]{0.5mm}{2.5mm}\;\Varid{e}\mskip1.5mu]} & \movesto & \ensuremath{\Conid{SX} [\mskip1.5mu \Varid{e}\mskip1.5mu]} \\
    \rulename{fail-r} & \ensuremath{\Conid{SX} [\mskip1.5mu \Varid{e}\;\rule[-0.3mm]{0.5mm}{2.5mm}\;\textbf{fail}\mskip1.5mu]} & \movesto & \ensuremath{\Conid{SX} [\mskip1.5mu \Varid{e}\mskip1.5mu]} \\
    \rulename{assoc-choice} & \ensuremath{\Conid{SX} [\mskip1.5mu (\Varid{e}_{\mathrm{1}}\;\rule[-0.3mm]{0.5mm}{2.5mm}\;\Varid{e}_{\mathrm{2}})\;\rule[-0.3mm]{0.5mm}{2.5mm}\;\Varid{e}_{\mathrm{3}}\mskip1.5mu]} & \movesto & \ensuremath{\Conid{SX} [\mskip1.5mu \Varid{e}_{\mathrm{1}}\;\rule[-0.3mm]{0.5mm}{2.5mm}\;(\Varid{e}_{\mathrm{2}}\;\rule[-0.3mm]{0.5mm}{2.5mm}\;\Varid{e}_{\mathrm{3}})\mskip1.5mu]} \\
    \rulename{choose} & \ensuremath{\Conid{SX} [\mskip1.5mu CX [\mskip1.5mu \Varid{e}_{\mathrm{1}}\;\rule[-0.3mm]{0.5mm}{2.5mm}\;\Varid{e}_{\mathrm{2}}\mskip1.5mu]\mskip1.5mu]} & \movesto & \ensuremath{\Conid{SX} [\mskip1.5mu CX [\mskip1.5mu \Varid{e}_{\mathrm{1}}\mskip1.5mu]\;\rule[-0.3mm]{0.5mm}{2.5mm}\;CX [\mskip1.5mu \Varid{e}_{\mathrm{2}}\mskip1.5mu]\mskip1.5mu]} & \text{if}~\ensuremath{CX} \ne \hole
    \end{array} \\

\iffullverse
    \\[-2mm]

    \textit{Effects} \\
    \begin{array}{lr@{\;}c@{\;}l}
    \rulename{succ-value} & \ensuremath{\textbf{succeeds}\,\{\Varid{v}\}} & \movesto & \ensuremath{\Varid{v}} \\
    \rulename{dec-fail} & \ensuremath{\textbf{decides}\,\{\textbf{fail}\}} & \movesto & \ensuremath{\textbf{fail}} \\
    \rulename{dec-value} & \ensuremath{\textbf{decides}\,\{\Varid{v}\}} & \movesto & \ensuremath{\Varid{v}} \\
    \end{array} \\
\fi
   \end{array}
   $$
  \caption{The Verse calculus: rewrite rules (2)}
    \label{fig:rules2}  \label{fig:choice}
\end{figure}
}

% ---------------------- Commenting out the wrong rules ------------------
% \begin{figure}
%   $$
%   \begin{array}{l}
%     \textbf{Syntax extension} \\
%     \begin{array}{lr@@{\;}c@@{\;}l}
%       \text{Expressions} & e  & ::= \cdots \vb |wrong| \\
%     \end{array} \\
%     \\
%     \textbf{Axiom extension} \\
%     \begin{array}{lr@@{\;}c@@{\;}ll}
%     \rulename{p-add-wrong} & |add^(hnf_1,hnf_2) | & \movesto & |wrong| & \text{if $|hnf_1| \neq k$ or $|hnf_2| \neq k$} \\
%     \rulename{p-gt-wrong} & |gt^(hnf_1,hnf_2) | & \movesto & |wrong| & \text{if $|hnf_1| \neq k$ or $|hnf_2| \neq k$} \\
%     \rulename{app-const-wrong} & |k^(e) | & \movesto & |wrong| \\
%     \rulename{unify-wrong-1} & |(lam x_1 e_1) = (lam x_2 e_2)| & \movesto & |wrong| \\
%     \rulename{unify-wrong-2} & |(lam x e) == hnf| & \movesto & |wrong|\\
%     \rulename{unify-wrong-3} & |hnf == (lam x e)| & \movesto & |wrong|\\
%     \rulename{occurs-check} & |x=V^[x]| & \movesto & |wrong| & \text{if }|V| \ne |hole|\\
% \iffullverse
%     \rulename{succeeds-fail} & |succeeds fail| & \movesto & |wrong| \\
%     \rulename{succeeds-choice-2} & |succeeds (v_1 choice v_2)| & \movesto & |wrong| \\
%     \rulename{succeeds-choice-n} & |succeeds (v_1 choice v_2 choice e)| & \movesto & |wrong| \\
%     \rulename{decides-choice-2} & |decides (v_1 choice v_2)| & \movesto & |wrong| \\
%     \rulename{decides-choice-n} & |decides (v_1 choice v_2 choice e)| & \movesto & |wrong| \\
% \fi
%     \rulename{wrong-seql} & |wrong; e| & \movesto & |wrong|  \\
%     \rulename{wrong-seqr} & |e; wrong| & \movesto & |wrong|  \\
%     \rulename{wrong-unifyl} & |wrong = e| & \movesto & |wrong|  \\
%     \rulename{wrong-unifyr} & |e = wrong| & \movesto & |wrong|  \\
%     \rulename{wrong-def} & |def x wrong| & \movesto & |wrong|  \\
%     \rulename{one-wrong} & |one wrong| & \movesto & |wrong| \\
%     \rulename{all-wrong} & |all (dots choice wrong choice dots)| & \movesto & |wrong| \\
%     \rulename{succeeds-wrong} & |succeeds wrong| & \movesto & |wrong| \\
%     \end{array}
%   \end{array}
%   $$
% \caption{The Verse calculus: axioms for |wrong|} \label{fig:wrong-axioms}
% \end{figure}
